% =====================================================================
\section{Part III: 25 Size and BE/ME Portfolios}
% =====================================================================

\paragraph{Sample.} 1926-07 to 2026-06, $T = 1200$, $N = 25$. Returns are in
excess of the one-month T-bill rate, in percent per month. Rows sort on size
(S = small, B = big); columns sort on BE/ME (L = low, H = high).

\subsection{Question (f), repeating Part I(a)}

\begin{table}[H]
\centering
\caption{Summary statistics, 25 size and BE/ME portfolios, percent per month.}
\small
\begin{tabular}{lrrrrrcrrrrr}
\toprule
& \multicolumn{5}{c}{Mean excess return} && \multicolumn{5}{c}{Standard deviation} \\
\cmidrule{2-6}\cmidrule{8-12}
& L & 2 & 3 & 4 & H && L & 2 & 3 & 4 & H \\
\midrule
S & 0.58 & 0.68 & 1.00 & 1.14 & 1.35 && 11.88 & 9.62 & 8.82 & 8.25 & 9.13 \\
2 & 0.66 & 0.95 & 0.98 & 1.04 & 1.24 &&  7.96 & 7.46 & 7.21 & 7.41 & 8.64 \\
3 & 0.73 & 0.91 & 0.92 & 1.01 & 1.11 &&  7.34 & 6.47 & 6.45 & 6.91 & 8.34 \\
4 & 0.73 & 0.79 & 0.86 & 0.97 & 1.03 &&  6.17 & 6.08 & 6.28 & 6.80 & 8.52 \\
B & 0.69 & 0.66 & 0.73 & 0.67 & 1.01 &&  5.34 & 5.24 & 5.56 & 6.50 & 8.48 \\
\midrule
& \multicolumn{5}{c}{Sharpe ratio} && \multicolumn{5}{c}{} \\
\cmidrule{2-6}
S & 0.049 & 0.071 & 0.113 & 0.138 & 0.148 &&&&&& \\
2 & 0.083 & 0.127 & 0.137 & 0.141 & 0.144 &&&&&& \\
3 & 0.100 & 0.141 & 0.142 & 0.145 & 0.133 &&&&&& \\
4 & 0.118 & 0.129 & 0.136 & 0.143 & 0.121 &&&&&& \\
B & 0.130 & 0.125 & 0.131 & 0.103 & 0.119 &&&&&& \\
\bottomrule
\end{tabular}
\label{tab:p3f_stats}
\vspace{2pt}\begin{minipage}{0.93\textwidth}\footnotesize
Monthly Sharpe ratio $= \bar R^e_i/\hat\sigma_i$. Market proxy (RM$-$RF):
mean $0.70$, standard deviation $5.31$, Sharpe ratio $0.131$.
\end{minipage}
\end{table}

\paragraph{Discussion.}
\begin{itemize}
\setlength{\itemsep}{4pt}\setlength{\parskip}{0pt}

\item \textbf{The dominant pattern is a value effect.} Average returns and
Sharpe ratios rise with BE/ME in every size quintile, and the spread is
steepest among small stocks.

\item \textbf{The size effect is conditional.} Small value beats big value,
but small growth earns less than big growth and is the worst of the 25
portfolios, with the lowest mean and the highest volatility.

\item \textbf{Unlike the industries, the dispersion is ordered.} The
industries in Part I span a similar range of Sharpe ratios with no ordering
to it. Here the dispersion lines up with the characteristic we sorted on,
which is exactly what a single beta will struggle to match.

\end{itemize}

\subsection{Question (f), repeating Part I(b)}

For each portfolio we estimate
\[
R_{it} - R_{ft} = \alpha_i + \beta_{iM}\,(R_{Mt} - R_{ft}) + e_{it},
\]
and test $H_0:\alpha_1 = \dots = \alpha_{25} = 0$ with the GRS statistic,
distributed $F(N,\,T-N-K) = F(25,\,1174)$ under the null.

\begin{table}[H]
\centering
\caption{CAPM time-series regressions, 25 size and BE/ME portfolios.
Intercepts in percent per month.}
\small
\begin{tabular}{lrrrrrcrrrrr}
\toprule
& \multicolumn{5}{c}{$\hat\alpha_i$} && \multicolumn{5}{c}{$t(\hat\alpha_i)$} \\
\cmidrule{2-6}\cmidrule{8-12}
& L & 2 & 3 & 4 & H && L & 2 & 3 & 4 & H \\
\midrule
S & $-0.53$ & $-0.29$ & 0.06 & 0.26 & 0.41 && $-2.20$ & $-1.59$ & 0.38 & 1.88 & 2.48 \\
2 & $-0.22$ & 0.09 & 0.15 & 0.20 & 0.29 && $-1.80$ & 0.88 & 1.53 & 1.88 & 2.12 \\
3 & $-0.13$ & 0.13 & 0.14 & 0.19 & 0.17 && $-1.43$ & 1.77 & 1.88 & 2.17 & 1.37 \\
4 & $-0.03$ & 0.03 & 0.09 & 0.17 & 0.07 && $-0.41$ & 0.50 & 1.36 & 1.96 & 0.57 \\
B & 0.02 & 0.00 & 0.06 & $-0.09$ & 0.10 && 0.51 & 0.01 & 0.90 & $-1.05$ & 0.72 \\
\midrule
& \multicolumn{5}{c}{$\hat\beta_{iM}$} && \multicolumn{5}{c}{} \\
\cmidrule{2-6}
S & 1.60 & 1.39 & 1.36 & 1.26 & 1.36 &&&&&& \\
2 & 1.27 & 1.23 & 1.20 & 1.21 & 1.37 &&&&&& \\
3 & 1.25 & 1.13 & 1.12 & 1.17 & 1.35 &&&&&& \\
4 & 1.09 & 1.09 & 1.10 & 1.16 & 1.39 &&&&&& \\
B & 0.96 & 0.94 & 0.96 & 1.09 & 1.30 &&&&&& \\
\bottomrule
\end{tabular}
\label{tab:p3f_capm}
\vspace{2pt}\begin{minipage}{0.93\textwidth}\footnotesize
GRS $= 3.37$, $p < 0.001$, $F(25, 1174)$. $\hat\Sigma$ uses denominator
$T-K-1$; the factor variance uses $T-1$.
\end{minipage}
\end{table}

\begin{itemize}
\setlength{\itemsep}{4pt}\setlength{\parskip}{0pt}

\item \textbf{The GRS test rejects.} $F = 3.37$ with $p < 0.001$: the market
proxy is not mean--variance efficient with respect to the 25 size and BE/ME
portfolios.

\end{itemize}

\subsection{Question (f), repeating Part I(d)}

\paragraph{Discussion.}
\begin{itemize}
\setlength{\itemsep}{4pt}\setlength{\parskip}{0pt}

\item \textbf{The intercepts are negative for growth and positive for
value.} The spread is widest among small stocks, from about $-0.5\%$ per
month for small growth to $+0.4\%$ for small value. Among the largest
stocks the intercepts are close to zero.

\item \textbf{The CAPM has most difficulty with small growth and small
value.} Small growth is overpriced and small value underpriced relative to
their betas, while big stocks are priced reasonably well. Only five
intercepts are significant on their own, but the residuals are strongly
correlated across portfolios, so the GRS test reads the alphas as one
systematic pattern rather than 25 small deviations.

\item \textbf{Beta cannot produce the spread, which points to a missing
factor for value and mispricing for small growth.} Within a size quintile
beta barely changes across BE/ME, and small growth has the highest beta of
all 25 portfolios despite the lowest return. The value spread moves as a
common component of the CAPM residuals, but it does not lose money in
market crashes, so it is not compensation for market risk. Small growth has
the highest skewness and idiosyncratic volatility of the 25 portfolios, a
negative alpha in every subsample, and stays mispriced even after adding a
value factor, which fits investors overpaying for lottery-like stocks.

\end{itemize}

\subsection{Question (g)}

We replace RM$-$RF with the in-sample tangency portfolio of the 25
portfolios, with weights
\[
w_T = \frac{\hat\Sigma^{-1}\hat\mu}{\mathbf{1}'\hat\Sigma^{-1}\hat\mu}
\]
estimated over the full sample, and rerun the 25 time-series regressions and
the GRS test.

\begin{table}[H]
\centering
\caption{GRS test with the in-sample tangency portfolio of the 25 size and
BE/ME portfolios as the market proxy, July 1926 -- June 2026.}
\small
\begin{tabular}{lrr}
\toprule
& RM$-$RF & In-sample tangency \\
\midrule
GRS $F$-statistic & 3.37 & 0.00 \\
$p$-value & $<0.001$ & 1.00 \\
Largest $|\hat\alpha_i|$ (\% per month) & 0.53 & 0.00 \\
Proxy Sharpe ratio (monthly) & 0.131 & 0.299 \\
\bottomrule
\end{tabular}
\label{tab:p3g}
\vspace{2pt}\begin{minipage}{0.93\textwidth}\footnotesize
Tangency weights sum to one, with gross exposure $13.86$ and individual
weights between $-1.44$ and $1.06$. All 25 intercepts are below $10^{-14}$.
\end{minipage}
\end{table}

\paragraph{Discussion.}
\begin{itemize}
\setlength{\itemsep}{4pt}\setlength{\parskip}{0pt}

\item \textbf{The GRS test fails to reject.} The statistic is zero and the
$p$-value is one: every intercept is zero up to rounding error, so the
in-sample tangency portfolio prices all 25 portfolios perfectly.

\item \textbf{The result is not surprising.} It is an identity. Because
$w_T \propto \hat\Sigma^{-1}\hat\mu$, each portfolio's covariance with the
tangency portfolio is proportional to its mean, so
$\hat\mu_i = \hat\beta_{iT}\,\hat\mu_T$ and
$\hat\alpha_i = \hat\mu_i - \hat\beta_{iT}\hat\mu_T = 0$ for every $i$,
whatever the data. In Sharpe-ratio terms, the proxy is already the
maximum-Sharpe-ratio combination of the test assets, so adding them cannot
improve it. The result therefore says nothing about the CAPM. The proxy is a
levered long--short portfolio of the test assets, not a market portfolio,
and it illustrates Roll's critique just as the perfect cross-sectional fit
did in Problem Set 3, Question (2).

\end{itemize}

\subsection{Question (h)}

Sample 1 contains odd-numbered months in even years and even-numbered months
in odd years; Sample 2 contains the rest, giving 600 months each. We estimate
tangency weights from Sample 1 moments and apply them to Sample 2 returns,
then reverse the roles. The combined, chronologically sorted series is the
out-of-sample tangency portfolio, which we use as the market proxy in the 25
time-series regressions.

\begin{table}[H]
\centering
\caption{GRS tests on the 25 size and BE/ME portfolios under three market
proxies, July 1926 -- June 2026.}
\small
\begin{tabular}{lrrr}
\toprule
& RM$-$RF & In-sample tangency & Out-of-sample tangency \\
\midrule
GRS $F$-statistic & 3.37 & 0.00 & 2.78 \\
$p$-value & $<0.001$ & 1.00 & $<0.001$ \\
Proxy Sharpe ratio (monthly) & 0.131 & 0.299 & 0.216 \\
\bottomrule
\end{tabular}
\label{tab:p3h}
\vspace{2pt}\begin{minipage}{0.93\textwidth}\footnotesize
Out-of-sample weights sum to one in each half, with gross exposure $14.81$
(estimated on Sample 1) and $18.44$ (Sample 2). The correlation between the
two weight vectors is $0.49$. All 25 intercepts are positive, between $0.14$
and $0.46\%$ per month.
\end{minipage}
\end{table}

\paragraph{Discussion.}
\begin{itemize}
\setlength{\itemsep}{4pt}\setlength{\parskip}{0pt}

\item \textbf{The GRS test rejects.} The out-of-sample tangency portfolio
leaves significant pricing errors in the 25 portfolios.

\item \textbf{It rejects because the weights are estimated with error.}
The identity in (g) holds only when the weights and the test use the same
sample moments. Here they do not, and the mean estimates are noisy enough
that the two halves produce quite different weights. A portfolio that is
efficient in its own half loses about a third of its Sharpe ratio in the
other, so it lies inside the frontier of the test assets. Its extra
volatility is mostly noise that shrinks every beta, so every portfolio's
average return exceeds what its beta implies and all 25 intercepts are
positive.

\item \textbf{The result is not surprising.} Optimized weights fit the
sample they come from and travel poorly, because mean returns are estimated
with large error. The portfolio still beats the market's Sharpe ratio and
lowers the GRS statistic, but the test asks whether the proxy is efficient,
not whether it beats the market, and out of sample it is not.

\end{itemize}

\paragraph{AI integration: Roll's critique.} Before comparing (g) and (h),
we asked Claude Sonnet 5.5 (chat mode) why the in-sample tangency portfolio
makes the GRS test fail to reject trivially, and why the out-of-sample
version does not. It answered that GRS is a monotone function of
$(SR^{*2} - SR_f^2)/(1 + SR_f^2)$, that an in-sample tangency portfolio
attains $SR_f = SR^*$ by construction, and that weights estimated on a
different sample leave the proxy inside the test-sample frontier, restoring
power. It added that estimation error in those weights makes the test
over-reject as a test of the underlying model.

\begin{itemize}
\setlength{\itemsep}{4pt}\setlength{\parskip}{0pt}

\item[(i)] \textbf{Yes, on both counts.} The in-sample tangency portfolio
produced a GRS statistic of zero and a $p$-value of one, exactly as the LLM
predicted.

\item[(ii)] \textbf{The LLM identified the exact identity.} It showed that
$\mathrm{Cov}(R_i, f) = \hat\mu_i$, so $\hat\beta_i = \hat\mu_i/\mu_f$ and
every $\hat\alpha_i$ is zero, rather than appealing loosely to overfitting.
It did overstate one point: it claimed the exact $F$ distribution applies
out of sample because the proxy is a given regressor. In our construction
each half's proxy uses weights estimated from the other half's test-asset
returns, so the proxy is not independent of the test sample and the $F$
distribution is only approximate.

\item[(iii)] \textbf{Yes, and the contrast reinforces Roll's critique.}
With the same assets, the same data and the same model, the verdict moves
from a perfect fit to a clear rejection depending only on how the proxy is
built. The LLM's own caveat makes the same point without naming Roll: the
out-of-sample rejection partly reflects estimation error in the proxy, so
neither result is a clean verdict on the CAPM. GRS can only ever test
whether the chosen proxy is efficient.

\end{itemize}

\subsection{Question (i)}

We replace RM$-$RF with the in-sample tangency portfolio of the 30
value-weight industry portfolios, estimated over the full sample, and rerun
the GRS test on the 25 size and BE/ME portfolios.

\begin{table}[H]
\centering
\caption{CAPM intercepts with the in-sample industry tangency portfolio as
the market proxy, percent per month.}
\small
\begin{tabular}{lrrrrrcrrrrr}
\toprule
& \multicolumn{5}{c}{$\hat\alpha_i$} && \multicolumn{5}{c}{$t(\hat\alpha_i)$} \\
\cmidrule{2-6}\cmidrule{8-12}
& L & 2 & 3 & 4 & H && L & 2 & 3 & 4 & H \\
\midrule
S & $-0.46$ & $-0.10$ & 0.26 & 0.47 & 0.62 && $-1.40$ & $-0.36$ & 1.06 & 2.03 & 2.44 \\
2 & $-0.17$ & 0.19 & 0.25 & 0.33 & 0.46 && $-0.81$ & 0.95 & 1.32 & 1.63 & 1.96 \\
3 & $-0.09$ & 0.14 & 0.19 & 0.27 & 0.31 && $-0.48$ & 0.86 & 1.14 & 1.46 & 1.37 \\
4 & $-0.10$ & 0.02 & 0.11 & 0.22 & 0.21 && $-0.64$ & 0.11 & 0.66 & 1.21 & 0.91 \\
B & $-0.12$ & $-0.07$ & 0.03 & $-0.05$ & 0.20 && $-1.00$ & $-0.54$ & 0.22 & $-0.31$ & 0.86 \\
\bottomrule
\end{tabular}
\label{tab:p3i}
\vspace{2pt}\begin{minipage}{0.93\textwidth}\footnotesize
GRS $= 3.51$, $p < 0.001$, against $3.37$ with RM$-$RF. Monthly Sharpe
ratio of the industry tangency portfolio: $0.236$.
\end{minipage}
\end{table}

\paragraph{Discussion.}
\begin{itemize}
\setlength{\itemsep}{4pt}\setlength{\parskip}{0pt}

\item \textbf{The industry tangency portfolio prices the 25 portfolios no
better than the market.} The GRS test still rejects, with a slightly larger
statistic than under RM$-$RF, and the intercepts keep the same pattern:
negative for small growth, positive and increasing in BE/ME elsewhere.

\item \textbf{This is not surprising.} The tangency portfolio is efficient
only relative to the assets it is built from. Industries are not sorted on
size or BE/ME, so their optimal combination carries little of the value
spread that drives the 25 portfolios' returns. A higher Sharpe ratio than
the market does not help: what matters is whether the proxy lies on the
frontier of the test assets, and the industry tangency portfolio points in a
different direction, as the low correlation between the two tangency
portfolios in (k) confirms.

\end{itemize}

\subsection{Question (j)}

We replace RM$-$RF with the in-sample tangency portfolio of the 10 past
return portfolios, estimated over the full sample, and rerun the GRS test on
the 25 size and BE/ME portfolios. Because the past return portfolios begin in
January 1927, this test uses 1927-01 to 2026-06 ($T = 1194$).

\begin{table}[H]
\centering
\caption{CAPM intercepts with the in-sample past-return tangency portfolio
as the market proxy, percent per month, 1927-01 to 2026-06.}
\small
\begin{tabular}{lrrrrrcrrrrr}
\toprule
& \multicolumn{5}{c}{$\hat\alpha_i$} && \multicolumn{5}{c}{$t(\hat\alpha_i)$} \\
\cmidrule{2-6}\cmidrule{8-12}
& L & 2 & 3 & 4 & H && L & 2 & 3 & 4 & H \\
\midrule
S & $-0.20$ & $-0.07$ & 0.35 & 0.50 & 0.74 && $-0.58$ & $-0.24$ & 1.39 & 2.10 & 2.77 \\
2 & $-0.17$ & 0.21 & 0.26 & 0.38 & 0.53 && $-0.77$ & 0.99 & 1.30 & 1.83 & 2.15 \\
3 & $-0.07$ & 0.12 & 0.19 & 0.31 & 0.39 && $-0.37$ & 0.71 & 1.09 & 1.60 & 1.65 \\
4 & $-0.11$ & 0.02 & 0.12 & 0.28 & 0.28 && $-0.66$ & 0.15 & 0.71 & 1.49 & 1.17 \\
B & $-0.07$ & $-0.07$ & 0.07 & $-0.00$ & 0.28 && $-0.54$ & $-0.52$ & 0.44 & $-0.02$ & 1.17 \\
\bottomrule
\end{tabular}
\label{tab:p3j}
\vspace{2pt}\begin{minipage}{0.93\textwidth}\footnotesize
GRS $= 3.36$, $p < 0.001$, $F(25, 1168)$, against $3.37$ with RM$-$RF over
the same sample. The tangency portfolio is mainly long Winner ($0.89$) and
short Loser ($-0.96$), with a monthly Sharpe ratio of $0.264$.
\end{minipage}
\end{table}

\paragraph{Discussion.}
\begin{itemize}
\setlength{\itemsep}{4pt}\setlength{\parskip}{0pt}

\item \textbf{How well do we price the 25 size and BE/ME portfolios?}
No better than the market. The GRS statistic is essentially unchanged from
the CAPM, but the intercepts now rise even more steeply with BE/ME, reaching
$0.74\%$ per month for small value, and their average absolute size
increases by about half.

\item \textbf{Is this surprising? Why or why not?} We expected the momentum
tangency portfolio to do better, since its Sharpe ratio is about twice the
market's. On reflection, the result is understandable. The test does not
reward a high Sharpe ratio; it asks whether the proxy is efficient relative
to the 25 portfolios, and a bet on winners over losers carries no
information about value. It even works against value: value portfolios move
with past losers, so the proxy expects little from them and their alphas
grow. The GRS statistic holds steady only because those larger alphas are
also noisier. A better portfolio is not necessarily a better proxy.

\end{itemize}
